\section{Related Work}
\label{sec:related}

Our study sits at the intersection of four threads: internal pre-planning in language models,
the effect of RL post-training on revision, behavioral sycophancy, and uncertainty
decomposition. We organize related work by these themes and position our contribution against
each.

\para{Internal pre-planning (supports C1).} A line of interpretability work shows that
decoder-only LMs encode global properties of their \emph{entire} upcoming response in early
hidden states, despite next-token training. \citet{dong2025emergent} train simple probes on
the prompt hidden state that predict a response's length, number of reasoning steps, final
multiple-choice answer, and eventual correctness---before the first token---and report a
``U-shaped'' commitment trajectory along generation. \citet{pal2023future} establish the
foundational primitive that future tokens are linearly decodable from a single hidden state,
and \citet{wu2024plan} formalize ``pre-caching'' versus ``breadcrumbs'' and show planning
grows with scale. \citet{reasoningstrength2025} show that the reasoning-token budget of a
reasoning model is linearly decodable and steerable, the most literal evidence for
``planning the end of the turn''---but, like the rest of this thread, their subjects are all
RL/distilled reasoning models with \emph{no base contrast}. \citet{bigelow2025planning} give
a causal, circuit-level definition of planning and find that instruction tuning \emph{selects
among} existing plans rather than creating planning, and that base models can hold
``competing plans.'' We reuse the prompt-position probing protocol of
\citet{dong2025emergent}, but extend the target to the model's \emph{own} answer and, crucially,
run it across a matched base-to-RL gradient---the comparison this thread omits.

\para{RL post-training and revision (C2, contested).} \citet{ufo2025} provide the strongest
direct support for the under-reaction reading: single-turn RLVR collapses revision, with RL
models repeating the identical answer across turns and the unique-answer ratio dropping after
RL across methods and scales. \citet{kumar2024score} similarly find ``behavior collapse,''
where naive RL locks in the first answer, and introduce the clean $\Delta^{i\rightarrow c}$ /
$\Delta^{c\rightarrow i}$ decomposition we adopt to separate beneficial fixes from harmful
flips. Yet the same work notes the base model can \emph{over}-react, making the direction
model-dependent. We connect this behavioral literature to an internal probe on matched models,
and find that---measured internally---it is the base model, not the RL model, that
pre-commits.

\para{Sycophancy and change-of-mind (C2, behavioral).} \citet{laban2023flipflop} run
challenger prompts (``are you sure?'') across ten models and report a 46\% flip rate and
$-17\%$ accuracy, with sycophancy \emph{universal} across models trained with and without
RLHF and \emph{not} monotonic in RLHF---a direct warning against a naive ``RL causes the
deficit'' reading. This over-reaction failure is the mirror image of pre-commitment; the two
coexist and are selected by model and task. Our two-task design (\arc and \gsm) reproduces
both ends of this axis within one family.

\para{Uncertainty decomposition (C3, open).} Methods exist for measuring and decomposing LLM
uncertainty into sources~\citep{uncertaintyprofiles2025} and for response-wise uncertainty
over reasoning steps, but no prior work rotates which part of a decomposable question is held
certain to test whether residual uncertainty redistributes. We introduce that protocol and
report the first measurements of it.

\para{Summary.} Prior work establishes that planning exists (C1) and that revision can
collapse \emph{or} over-fire after training (C2), but never on a matched base-vs-RL probe
contrast, and never tests rotating certainty (C3). We close all three gaps in one study.
